\section{Mapeo general x86 $\to$ RISC-V}
\label{sec:mapeo}

La siguiente tabla resume las equivalencias que se irán desarrollando en las secciones siguientes. Varios elementos del
TP3 no tienen contraparte directa y su función la cumple otro mecanismo; se indica con ``---''.

\begin{center}
\small
\begin{tabular}{p{6.2cm} | p{8.6cm}}
\textbf{TP3 x86} & \textbf{RISC-V (este port)} \\
\hline
GDT + segmentación (code/data, nivel 0/3) & --- (no hay segmentos). Permisos por PTE Sv39 (\texttt{U/R/W/X}). \\
Modo real $\to$ modo protegido + A20 & M-mode $\to$ S-mode vía \texttt{mret}; PMP abierta. \\
CR0.PG + CR3 & CSR \texttt{satp} con modo SV39 + ASID + PPN. \\
IDT + 255 ISR stubs & CSR \texttt{stvec} (único) + despacho en C por \texttt{scause}. \\
Stack separada en cambio de nivel (SS0:ESP0 en TSS) & Stack dedicada de trap (\texttt{trap\_stack\_top}) y salto de
\texttt{gp} al del kernel. \\
\texttt{pushad}, selectores + \texttt{iret} & Guardado manual en \texttt{TrapFrame} (GPRs +
\texttt{sepc}/\texttt{sstatus}/\texttt{scause}/\texttt{stval}) y \texttt{sret}. \\
PIT 8254 + IRQ0 + PIC \texttt{pic\_finish1} & CLINT \texttt{mtimecmp}/\texttt{mtime} en M-mode $\to$ SSIP; \texttt{csrc
sip, SSIP}. \\
Teclado PS/2 IRQ1 + scancodes & UART 16550 por polling; teclas ASCII. \\
VGA texto en \texttt{0xB8000} & UART + secuencias ANSI (80$\times$50 emulado). \\
TSS + cambio de tarea por hardware (\texttt{jmp <sel>:<off>}) & \texttt{TrapFrame} por tarea + \texttt{sret} (cambio por
software). \\
\texttt{int 0x2F} (syscall) & \texttt{ecall}: \texttt{a7}=id, \texttt{a0}=arg/retorno. \\
\texttt{invlpg} / recargar CR3 & \texttt{sfence.vma} (por VA o global, opcionalmente por ASID). \\
\texttt{bochsrc} / Bochs & QEMU \texttt{virt} con \texttt{-bios none}. \\
\end{tabular}
\end{center}

Las referencias a instrucciones y CSRs siguen la convención de~\cite{riscv-unpriv} (ISA no privilegiada)
y~\cite{riscv-priv} (ISA privilegiada, CSRs, Sv39).
